\documentstyle[% 


righttag% 


,11pt% 


,amssymb% 


,russian%               remove in Eng. 


,izvru%                 Set izven% in Eng. 


]{amsart} 


 


%  


 


\newcommand{\dst}{\displaystyle} 


\newcommand{\fuckprime}{\overset{\scriptstyle\prime}{\phantom{i}}{}\!\!\!} 


\newcommand{\ord}{\operatorname{ord}} 


\newcommand{\tts}[1]{} 


\renewcommand{\baselinestretch}{0.96}


 


\theoremstyle{plain} 


\theorembodyfont{\it} 


\newtheorem{lemnonum}{\hskip\parindent Лемма} 


\renewcommand{\thelemnonum}{} 


 


%\hoffset=-15mm%                  For Eng. 


\setcounter{page}{20} 


\god{2000} 


\nomer{1~(452)} %                        Remove in Eng. 


%\nomer{Vol.\,44, No.\,1}%               Set in Eng. 


%\str{\pageref{begin}--\pageref{end}}%  Set in Eng. 


\udk{517.948\,:\,517.544} 


 


\author{\it Н.А.\,Дегтяренко} 


% NB:   ^^ remove in Eng. 


\title{Решение в замкнутой форме 


интегрального уравнения типа свертки


в гиперэллиптическом случае} 


 


\date{} 


% \pagestyle{myheadings}%         Set in Eng. 


\pagestyle{plain} %              Remove in Eng. 


\begin{document} 


% \thispagestyle{plain}%          Set in Eng. 


\maketitle 


\label{begin} 


 


{\bf 1.} Будем рассматривать интегральное уравнение типа свертки вида 


\begin{multline} 


\lambda(x)\varphi(x)+\frac{1}{\sqrt{2\pi}} 


\int^\infty_0 k_1(x-t)\varphi(t)dt+\frac{1}{\sqrt{2\pi}} 


\int^0_{-\infty}k_2(x-t)\varphi(t)dt+\mu(x)\ov{\varphi(x)}+ \\ 


% 


+\frac{1}{\sqrt{2\pi}}\int^\infty_0 k_3(x-t) 


\ov{\varphi(t)}dt+\frac{1}{\sqrt{2\pi}}\int^0_{-\infty} 


k_4(x-t)\ov{\varphi(t)}dt=f(x),\quad -\infty<x<\infty, 


\end{multline} 


где $\lambda(x)$ и $\mu(x)$ --- вещественные кусочно-постоянные функции


$$ 


\lambda(x)= 


\begin{cases} 


\lambda_1, & x>0; \\ 


\lambda_2, & x<0, 


\end{cases} 


\qquad 


\mu(x)= 


\begin{cases} 


\lambda_3, & x>0;\\ 


\lambda_4, & x<0. 


\end{cases} 


$$ 


Это уравнение впервые было рассмотрено в ([1], c.\,33--34). 


 


Искомая функция $\varphi(x)$, свободный член $f(x)$ и разностные ядра 


считаются 


принадлежащими классу $\{0\}$. Классы $\{0\}$ и $\{\{0\}\}$ были 


введены в ([2], с.\,12). 


При помощи аппарата преобразования Фурье уравнение (1) приводится в 


([1], c.\,40--41) к краевому условию вида 


\begin{multline} 


[\lambda_1+K_1(x)]\Phi^+(x)-[\lambda_2+K_2(x)]\Phi^-(x)+ 


[\lambda_3+K_3(x)]\ov{\Phi^+(-x)}- \\


%


-[\lambda_4+K_4(x)]\ov{\Phi^-(-x)}=F(x),\quad 


-\infty<x<\infty, 


\end{multline} 


где $\Phi^\pm(x)\in\{\{0\}\}$ и являются предельными значениями на 


действительной оси 


аналитических соответственно в верхней и нижней полуплоскостях 


функций $\Phi^\pm(z)$. Присоединяя к (2) равенство, полученное из него 


комплексным 


сопряжением и заменой $x$ на $(-x)$, и исключая из них $\ov{\Phi^+(-x)}$, 


получим краевое 


условие вида 


\begin{equation} 


\Delta_1(x)\Phi^+(x)=C(x)\Phi^-(x)+D(x)\ov{\Phi^-(-x)}+ 


G(x),\quad -\infty<x<\infty, 


\end{equation} 


где введены обозначения


\begin{align*} 


C(x)&=[\lambda_2+K_2(x)][\lambda_1+\ov{K_1(-x)}]


-[\lambda_3+K_3(x)][\lambda_4+\ov{K_4(-x)}];\\ 


% 


D(x)&=[\lambda_1+\ov{K_1(-x)}][\lambda_4+K_4(x)]- 


[\lambda_3+K_3(x)][\lambda_2+\ov{K_2(-x)}];\\ 


% 


G(x)&=[\lambda_1+\ov{K_1(-x)}]F(x)-[\lambda_3+K_3(x)] 


\ov{F(-x)};\\ 


% 


\Delta_1(x)&=[\lambda_1+K_1(x)][\lambda_1+\ov{K_1(-x)}]- 


[\lambda_3+K_3(x)][\lambda_3+\ov{K_3(-x)}]. 


\end{align*} 


Функции $\Delta_1(x)$, $C(x)$, $D(x)$, $G(x)$ принадлежат классу 


$\{\{0\}\}$. Обозначения для 


коэффициентов задачи (3) взяты из ([1], c.\,41). 


Полученная краевая задача (3) и уравнение (1) эквивалентны в том 


смысле, что, найдя решение $\Phi^\pm(x)$ краевой задачи, решение 


уравнения получим 


по формуле 


$$ 


\varphi(x)=\frac{1}{\sqrt{2\pi}}\int^\infty_{-\infty} 


[\Phi^+(t)-\Phi^-(t)]e^{-itx}dt. 


$$ 


При этом линейно независимым решениям краевой задачи (3) 


соответствуют линейно независимые решения уравнения (1) и наоборот. 


\medskip 


 


{\bf 2.} Зафиксируем целое неотрицательное число $h$. Обозначим через 


$r_1,r_2,\dotsc,r_{h+1}$ попарно 


различные фиксированные отрицательные вещественные числа, а через 


$r_{2h+2},r_{2h+1},\dotsc,r_{h+2}$ --- числа, соответственно 


симметричные им относительно нуля. 


Введем обозначение $U:=\bigcup\limits^h_{i=0}(r_{2i+1},r_{2i+2})$. 


Будем 


предполагать, что коэффициенты краевого условия (3) удовлетворяют 


следующим 


тождествам: 


\begin{equation} 


\begin{split} 


C(x)&\equiv 0\quad\text{при }\;x\in U,\\ 


D(x)&\equiv 0\quad\text{при }\; x\in R\setminus\ov U. 


\end{split} 


\end{equation} 


Тогда краевое условие (3) запишется в виде


\begin{equation} 


\begin{split} 


\Delta_1(x)\Phi^+(x)&=D(x)\ov{\Phi^-(-x)}+G(x)\quad 


\text{при }\; x\in U,\\ 


\Delta_1(x)\Phi^+(x)&=C(x)\Phi^-(x)+G(x)\phantom{-}\quad 


\text{при }\;x\in R\setminus\ov U. 


\end{split} 


\end{equation} 


В силу принадлежности функций $C(x)$ и $D(x)$ классу $\{\{0\}\}$ 


выполняются следующие 


равенства: 


\begin{equation} 


C(r_i)=0,\ \; D(r_i)=0\ \; \text{при}\ \; i=\ov{1,2h+2}. 


\end{equation} 


Будем предполагать, что также выполняются условия 


\begin{equation} 


\lambda_i+K_i(r_j)\ne 0,\ \; i=1,2;\ \; j=\ov{1,2h+2}. 


\end{equation} 


Тогда из (6) легко получить 


\begin{equation} 


\lambda_i+K_i(r_j)\ne 0,\ \; i=3,4;\ \; j=\ov{1,2h+2}. 


\tag{7$'$} 


\end{equation} 


 


\begin{lemnonum} 


Коэффициент $\Delta_1(x)$ в точках $x=r_j$, $j=\ov{1,2h+2}$, обращается 


в нуль. 


\end{lemnonum} 


 


\begin{pf} 


Рассмотрим любую точку $x=r_j$, $j=\ov{1,2h+2}$. Из условий (6) 


следуют равенства 


\begin{align*} 


[\lambda_2+\ov{K_2(-r_j)}] 


[\lambda_1+K_1(r_j)]&=[\lambda_3+\ov{K_3(-r_j)}] 


[\lambda_4+K_4(r_j)];\\ 


% 


[\lambda_2+K_2(r_j)][\lambda_1+\ov{K_1(-r_j)}]&= 


[\lambda_3+K_3(r_j)][\lambda_4+\ov{K_4(-r_j)}]; \\ 


% 


[\lambda_1+\ov{K_1(-r_j)}][\lambda_4+K_4(r_j)]&= 


[\lambda_3+K_3(r_j)][\lambda_2+\ov{K_2(-r_j)}];\\ 


% 


[\lambda_1+K_1(r_j)][\lambda_4+\ov{K_4(-r_j)}]&= 


[\lambda_3+\ov{K_3(-r_j)}] 


[\lambda_2+K_2(r_j)]. 


\end{align*} 


Перемножим все четыре равенства между собой и сократим общие множители. 


Выпишем результат 


$[\lambda_1+K_1(r_j)]^2[\lambda_1+\ov{K_1(-r_j)}]^2= 


[\lambda_3+K_3(r_j)]^2[\lambda_3+\ov{K_3(-r_j)}]^2$. 


Из этого равенства следуют две возможности: 


 


1) $[\lambda_1+K_1(r_j)][\lambda_1+\ov{K_1(-r_j)}]= 


[\lambda_3+K_3(r_j)][\lambda_3+\ov{K_3(-r_j)}]$, 


т.\,е. $\Delta_1(r_j)=\Delta_1(-r_j)=0$, 


 


2) $[\lambda_1+K_1(r_j)][\lambda_1+\ov{K_1(-r_j)}]= 


-[\lambda_3+K_3(r_j)][\lambda_3+\ov{K_3(-r_j)}]$. 


 


\noindent Покажем, что вторая возможность не реализуется. Предположим 


противное. Тогда 


выразим $C(-r_j)$ следующим образом:


\begin{multline*} 


C(-r_j)=-[\lambda_2+K_2(-r_j)]


\frac{[\lambda_3+\ov{K_3(r_j)}][\lambda_3+K_3(-r_j)]} 


{[\lambda_1+K_1(-r_j)]}-[\lambda_3+K_3(-r_j)] 


[\lambda_4+\ov{K_4(r_j)}] = \\ 


% 


=-\frac{[\lambda_3+K_3(-r_j)]}{[\lambda_1+K_1(-r_j)]} 


(\ov{D(r_j)}+2[\lambda_2+K_2(-r_j)] 


[\lambda_3+\ov{K_3(r_j)}])= \\


%


=-\frac{2[\lambda_3+K_3(-r_j)][\lambda_2+K_2(-r_j)] 


[\lambda_3+\ov{K_3(r_j)}]} 


{[\lambda_1+K_1(-r_j)]}=0. 


\end{multline*} 


Следовательно, хотя бы один из множителей числителя равен нулю. 


Получили противоречие с условиями (7). Полученное противоречие 


доказывает, что второе равенство не выполняется. Следовательно, имеет 


место только первое равенство для любой точки $x=r_j$, $j=\ov{1,2h+2}$.  


\end{pf}


 


Будем решать задачу (5), предполагая, что $\Delta_1(x)$, $C(x)$, $D(x)$ 


в соответствующих 


окрестностях точек $x=r_j$, $j=\ov{1,2h+2}$, (т.\,е. там, где 


коэффициенты отличны от 


тождественного нуля) являются бесконечно малыми функциями одного и 


того же порядка. Предположим также, что свободный член $G(x)$ в 


окрестностях этих точек имеет порядок не ниже, чем $\Delta_1(x)$. 


Запишем эти 


условия в виде следующих соотношений:


\begin{equation} 


\begin{split} 


C(x)&=O^*(\Delta_1(x))\quad\text{при}\ \; 


\begin{cases} 


x\to r_i,\ x<r_i, & i=2k+1,\ k=\ov{0,h};\\ 


x\to r_i,\ x>r_i, & i=2k+2,\ k=\ov{0,h}, 


\end{cases} 


\\ 


D(x)&=O^*(\Delta_1(x))\quad\text{при}\ \; 


\begin{cases} 


x\to r_i,\ x>r_i, & i=2k+1,\ k=\ov{0,h};\\ 


x\to r_i,\ x<r_i, & i=2k+2,\ k=\ov{0,h}, 


\end{cases} 


\\


G(x)&=O(\Delta_1(x))\phantom{^*}\quad\text{при}\quad 


x\to r_i,\ \; i=\ov{1,2h+2}. 


\end{split} 


\end{equation} 


 


В силу принадлежности классу $\{\{0\}\}$ пределы на бесконечности 


коэффициентов $\Delta_1(x)$, $C(x)$ 


задачи (5), свободного члена $G(x)$ и искомых функций $\Phi^\pm(x)$ 


равны нулю. Будем 


предполагать, что в окрестности бесконечно удаленной точки 


$\Delta_1(x)$ и $C(x)$ 


являются бесконечно малыми функциями одного и того же порядка, а $G(x)$ 


имеет в окрестности этой точки порядок не ниже, чем $\frac{C(x)}{x}$, 


т.\,е. 


\begin{equation} 


\begin{split} 


\Delta_1(x)&=O^*(C(x))\quad\text{при}\ \; x\to\infty,\\ 


G(x)&=O\Big(\frac{C(x)}{x}\Big)\quad\text{при}\ \; x\to\infty. 


\end{split} 


\end{equation} 


Сделаем следующие предположения: 


\begin{equation} 


\begin{split} 


\forall\,x\in R\setminus\{r_1,r_2,\dotsc,r_{2h+2}\}: 


\Delta_1(x)\ne 0,\\ 


\forall\,x\in R\setminus\ov U:C(x)\ne 0,\\ 


\forall\,x\in U:D(x)\ne 0. 


\end{split} 


\end{equation} 


Будем решать задачу, предполагая , что выполняются условия (4), 


(7)--(10). Из этих ограничений следует, что если решение задачи (5) 


существует, то оно будет принадлежать классу ограниченных функций, 


имеющих нуль кратности, большей либо равной единице лишь в одной точке 


(бесконечно удаленной). Заметим, что при этом предельные значения 


сверху и снизу $\Phi^\pm(x)$ на вещественной оси функций $\Phi^\pm(z)$ 


будут принадлежать классу $\{\{0\}\}$. 


\medskip 


 


{\bf 3.} Введем вектор-функцию 


\begin{equation} 


\begin{pmatrix} 


\Phi_1(z) \\ \Phi_2(z) 


\end{pmatrix} 


=\bigg(\frac{\Phi(z)}{\Phi(-\ov z)}\bigg). 


\end{equation} 


Очевидно, что функции $\Phi_1(z)$ и $\Phi_2(z)$ будут связаны между 


собой тождеством 


\begin{equation} 


\Phi_1(z)\equiv\ov{\Phi_2(-\ov z)}. 


\end{equation} 


Отсюда следует, что при $x\in\Bbb R$ выполняются следующие тождества:


\begin{equation} 


\Phi^+_1(x)\equiv\ov{\Phi^+_2(-x)},\quad 


\Phi^-_1(x)\equiv\ov{\Phi^-_2(-x)}. 


\end{equation} 


Здесь $\Phi^\pm_1(x)$, $\Phi^\pm_2(x)$ --- предельные значения сверху 


$(+)$ и снизу $(-)$ на действительной 


оси функций $\Phi_1(z)$, $\Phi_2(z)$. Присоединим к каждому равенству 


(5) равенство, полученное 


из него комплексным сопряжением и заменой $x$ на $(-x)$. С учетом формул 


(11)--(13) получим краевое условие вида 


\begin{equation} 


\begin{split} 


\Delta_1(x)\Phi^+_1(x)&=D(x)\Phi^-_2(x)+G(x), \\ 


\ov{\Delta_1(-x)}\Phi^+_2(x)&= 


\ov{D(-x)}\Phi^-_1(x)+\ov{G(-x)},\quad 


\text{\raisebox{5mm}[1pt]{при $x\in U;$}} \\ 


% 


\Delta_1(x)\Phi^+_1(x)&=C(x)\Phi^-_1(x)+G(x), \\ 


\ov{\Delta_1(-x)}\Phi^+_2(x)&= 


\ov{C(-x)}\Phi^-_2(x)+\ov{G(-x)},\quad 


\text{\raisebox{5mm}[1pt]{при $x\in R\setminus\ov U.$}} 


\end{split} 


\end{equation} 


Полученная система уравнений (14) эквивалентна системе уравнений (5). 


Заметим, что если в уравнениях (14) заменить $x$ на $(-x)$ и все функции 


заменить комплексно-сопряженными, то условия сопряжения (14) перейдут 


сами в себя. Система (14) есть задача Римана для двух пар функций с 


линией разрыва $R$. Cформулируем ее более подробно. 


{\it  


Найти все функции $\Phi_1(z)$, $\Phi_2(z)$, аналитические в 


$C\setminus R$, предельные значения которых на 


действительной оси принадлежат классу $\{\{0\}\}$ и удовлетворяют 


равенствам $(14)$, предполагая, что на коэффициенты краевого условия 


накладываются ограничения вида $(4)$, $(7)$--$(10)$. 


} 


%\smallskip


 


{\bf 4.} Задачу (14) будем решать сведением к скалярной задаче Римана 


на римановой 


поверхности рода $h$, задаваемой уравнением 


\begin{equation} 


w^2=\prod^{2h+2}_{k=1}(z-r_k). 


\end{equation} 


Рассмотрим расширенную комплексную плоскость с разрезами от $r_{2i+1}$ 


до $r_{2i+2}$ ($i=\ov{0,h}$) по 


вещественной оси. Обозначим через $D$ полученную область, а через 


$D\cup\partial D$ --- ее 


компактификацию Мазуркевича (область с краем) ([3], c.\,24). На 


разрезанной плоскости выделим две однозначные непрерывные ветви корня 


$w(z)=\sqrt{\prod\limits^{2h+2}_{k=1}(z-r_k)}$. Через 


$w(z)=\sqrt{\prod\limits^{2h+2}_{k=1}(z-r_k)}$ обозначим ту 


непрерывную ветвь, которая при $z\to\infty$ имеет разложение


$\sqrt{\prod\limits^{2h+2}_{k=1}(z-r_k)}=z^{h+1}+\dotsb$.


Возьмем два экземпляра $D\cup\partial D$, которые обозначим через 


$D^+\cup\partial D^+$ и $D^-\cup\partial D^-$. 


Расположив $D^+\cup\partial D^+$ над $D^-\cup\partial D^-$, 


склеим их ``накрест'' вдоль разрезов $[r_{2i+1},r_{2i+2}]$ 


($i=\ov{0,h}$). В результате получим риманову 


поверхность $\frak R$ рода $h$, алгебраическое уравнение которой 


совпадает с 


(15). Будем считать, что 


$$ 


D^\pm=\{(z,\ \pm w(\Dot z))\mid z\in D\}. 


$$ 


Каждый лист поверхности $\frak R$ имеет бесконечно удаленную точку. 


Обозначим 


эти точки соответственно символами $(\infty,\infty_1)$, 


$(\infty,\infty_2)$. За локальный параметр окрестностей 


вышеуказанных точек принимается $t=\frac{1}{z}$. Точки вида $(r_k,0)$ 


($k=\ov{1,2h+2}$) являются единственными 


точками ветвления поверхности $\frak R$. Переменная $t=\sqrt{z-r_k}$ 


используется в качестве 


локального параметра окрестностей точек ветвления. В регулярных 


точках $(z,w)\in\frak R$ за локальный параметр принимается $z$. 


 


Краю $\partial D$ припишем ориентацию, указанную на рис.\,1. 


\begin{center} 


\begin{picture}(400,140)(0,-15) 


\put(30,120){\special{em:graph degt.pcx}} 


\put(180,-10){\mbox Рис.~1.} 


\end{picture} 


\end{center}


Рисунок 1 соответствует случаю $h=3$. 


Область $D^+$ остается слева при движении 


вдоль края в положительном направлении. Условимся считать, что 


интервалы $(-\infty;r_1)$, $(r_{2i},r_{2i+1})$ ($i=\ov{1,h}$); 


$(r_{2h+2};+\infty)$ 


вещественной оси на листе $D^+$ ориентированы слева направо, а на 


листе $D^-$ --- в противоположном направлении (на рис.\,1 пунктиром 


обозначается линия на нижнем листе). Канонические сечения $a_i$ и $b_i$ 


($i=\ov{1,h}$) 


на $\frak R$ выбираем и ориентируем стандартным образом ([4], 


cс.\,119, 155). 


 


На $\frak R$ введем новую неизвестную кусочно-аналитическую функцию 


$F(z,w)$, $(z,w)\in\frak R$, полагая 


\begin{equation} 


F(z,w)= 


\begin{cases} 


\Phi_1(z), & (z,w)\in D^+; \\ 


\Phi_2(z), & (z,w)\in D^-. 


\end{cases} 


\end{equation} 


Из (12) следует, что она должна удовлетворять условию 


\begin{equation} 


F(z,w)=\ov{F(-\ov z,\,-w(\ov z))}. 


\end{equation} 


Функция $F(z,w)$ является разрывной вдоль общей границы областей $D^+$ 


и $D^-$. 


Обозначим эту границу $L_1$. Через $L_2$ обозначим объединение множеств 


вида $\{(x,\vartheta)\mid x\in(-\infty;r_1)\}$, 


$\{(x,\vartheta)\mid x\in(r_{2i},r_{2i+1}),\ i=\ov{1,h}\}$, 


$\{(x,\vartheta)\mid x\in(r_{2h+2};+\infty)\}$, 


где $\vartheta^2=\prod\limits^{2h+2}_{k=1}(x-r_k)$, 


$(x,\vartheta)\in\frak R$. 


 


Выпишем краевое условие на контуре $L:=L_1\cup L_2$ для функции 


$F(z,w)$. 


Введем обозначения 


\begin{align*} 


\wt\Delta_1(x,\vartheta)&= 


\begin{cases} 


D(x)/\Delta_1(x) &\text{на верхних берегах всех разрезов}; \\ 


\ov{\Delta_1(-x)}/\ov{D(-x)}& 


\text{на нижних берегах всех разрезов}, 


\end{cases} 


\\ 


\wt G_1(x,\vartheta)&= 


\begin{cases} 


G(x)/\Delta_1(x) &\text{на верхних берегах всех разрезов};\\ 


\ov{-G(-x)}/\ov{D(-x)}& 


\text{на нижних берегах всех разрезов}, 


\end{cases} 


\\ 


\wt\Delta_2(x,\vartheta)&= 


\begin{cases} 


C(x)/\Delta_1(x) & \text{при }\; (x,\vartheta)\in L_2, 


\ \; \vartheta>0; \\ 


\ov{\Delta_1(-x)}/\ov{C(-x)}& \text{при }\; 


(x,\vartheta)\in L_2,\ \; \vartheta<0, 


\end{cases} 


\\ 


\wt G_2(x,\vartheta)&= 


\begin{cases} 


G(x)/\Delta_1(x) & \text{при }\; (x,\vartheta)\in L_2, 


\ \; \vartheta>0; \\ 


\ov{-G(-x)}/\ov{C(-x)} & \text{при }\; (x,\vartheta)\in L_2, 


\ \; \vartheta<0, 


\end{cases} 


\\ 


\Delta(x,\vartheta)&= 


\begin{cases} 


\wt\Delta_1(x,\vartheta) & \text{при }\;(x,\vartheta)\in L_1; \\ 


\wt\Delta_2(x,\vartheta) & \text{при }\;(x,\vartheta)\in L_2, 


\end{cases} 


\\ 


G_1(x,\vartheta)&= 


\begin{cases} 


\wt G_1(x,\vartheta) & \text{при }\;(x,\vartheta)\in L_1; \\ 


\wt G_2(x,\vartheta) & \text{при }\;(x,\vartheta)\in L_2. 


\end{cases} 


\end{align*} 


Учитывая эти обозначения и (16), получим краевое условие на $\frak R$ 


\begin{equation} 


F^+(x,\vartheta)=\Delta(x,\vartheta)F^-(x,\vartheta)+ 


G_1(x,\vartheta),\quad (x,\vartheta)\in L, 


\end{equation} 


где $F^\pm(x,\vartheta)$ --- предельные значения слева и справа функции 


$F(z,w)$ на $L$. 


 


Заметим, что если в краевом условии (18) заменить $(x,\vartheta(x))$ на 


$(-x,-\vartheta(x))$ и заменить все 


функции комплексно-сопряженными, то оно при выполнении условия (17)


перейдет само в себя, т.\,к. коэффициент задачи $\Delta(x,\vartheta)$ и 


свободный член $G_1(x,\vartheta)$ 


удовлетворяют условиям 


\begin{equation} 


\Delta(x,\vartheta(x))=\frac{1}{\ov{\Delta(-x,-\vartheta(x))}},\quad 


G_1(x,\vartheta(x))=\ov{-G_1(-x,-\vartheta(x))} 


\Delta(x,\vartheta(x)). 


\end{equation} 


Эти условия проверяются непосредственно. С другой стороны, если 


$F(z,w(z))$ 


является решением краевой задачи (18), то $\ov{F(-\ov z,-w(\ov z))}$ 


также является решением 


этой задачи в силу (19). Таким образом, для того чтобы общее решение 


краевой 


задачи (18) давало общее решение задачи (14), необходимо и 


достаточно, чтобы выполнялось условие симметрии (17). Поскольку 


решение и коэффициенты задачи (18) принадлежат г\"ельдеровскому классу 


($H$-непрерывны), то для ее решения можно воспользоваться теоремами и 


формулами из [4]. Из постановки задачи (14) следует, что дивизор $J$ 


([4], c.\,129), определяющий класс функций, которому принадлежат решение 


(если оно существует) и свободный член задачи (18), имеет вид 


$J=(\infty,\infty_1)^{-1}(\infty,\infty_2)^{-1}$. 


\medskip 


 


{\bf 5.} Рассмотрим однородную задачу 


\begin{equation} 


\wh F^+(x,\vartheta)=\Delta(x,\vartheta)\wh F^-(x,\vartheta), 


\quad (x,\vartheta)\in L. 


\end{equation} 


Соответствующая союзная задача для дифференциалов имеет вид 


\begin{equation} 


d\Psi^+(x,\vartheta)=\frac{1}{\Delta(x,\vartheta)} 


d\Psi^-(x,\vartheta),\quad (x,\vartheta)\in L. 


\end{equation} 


Сначала будем решать задачи (20) и (21) без учета условия симметрии, а 


затем 


обеспечим его выполнение. 


 


Множество точек контура $L$ обозначим символом $\{L\}$ и будем считать 


его 


топологическим пространством с индуцированной топологией. Множество 


$\{L\}\setminus\{(r_1,0),(r_2,0),\dotsc,(r_{2h+2},0),\linebreak


(\infty,\infty_1),(\infty,\infty_2)\}$ 


распадается на конечное число связных компонент $\{L_j\}$, каждая из 


которых 


есть гладкая открытая жорданова кривая, гомеоморфная интервалу $(0;1)$ 


числовой оси. Кривые $L_j$ назовем {\it кривыми контура} $L$. Каждая 


кривая $L_j$ имеет 


начальную и концевую точки, которые принадлежат множеству 


$\{(r_1,0),(r_2,0),\dotsc,(r_{2h+2},0),


(\infty,\infty_1),(\infty,\infty_2)\}$. 


На каждой 


из них задана $H$-непрерывная функция $\Delta_j(x,\vartheta)$, конечная 


и не обращающаяся в нуль, 


$H$-непрерывно продолжимая на начальную и концевую точки кривой $L_j$, 


причем предельные значения являются конечными и отличными от нуля в 


силу (8)--(10). На каждой кривой $L_j$ можно выделить однозначную 


ветвь логарифма $\ln\Delta_j(x,\vartheta)$. Выделим произвольным образом 


указанные ветви, 


зафиксируем их и будем использовать в дальнейших формулах. 


 


Для решения возникших задач удобно использовать формулы и теоремы из 


[4]. Чтобы решить задачи (18), (20), (21), необходимо указать 


комплексно-нормированный базис абелевых дифференциалов первого рода 


на $\frak R$. Нормируя базис абелевых дифференциалов первого рода 


поверхности $\frak R$\; 


$\frac{dz}{w},\frac{z\,dz}{w},\dotsc,\frac{z^{h-1}dz}{w}$, 


$(z,w)\in\frak R$,


относительно выбранных канонических сечений, получим 


комплексно-нормированный базис 


$du_1(z,w),du_2(z,w),\dotsc,du_h(z,w)$. 


Дифференциалы этого базиса вычисляются по формулам ([4], c.\,155) 


$$ 


du_j(z,w)= 


\frac{ 


\begin{vmatrix} 


0 & \dst\frac{dz}{w} & \dotsc & \dst\frac{z^{h-1}dz}{w} \\[3mm] 


-\delta_{j1} & \dst \int_{a_1}\dst\frac{d\tau}{\xi} & \dotsc & 


\dst\int_{a_1}\frac{\tau^{h-1}d\tau}{\xi} \\ 


\dotsc & \dotsc & \dotsc & \dotsc \\ 


-\delta_{jh} & \dst\int_{a_h}\dst\frac{d\tau}{\xi} & \dotsc & 


\dst\int_{a_h}\dst\frac{\tau^{h-1}d\tau}{\xi} 


\end{vmatrix} 


} 


{ 


\begin{vmatrix} 


\dst\int_{a_1}\dst\frac{d\tau}{\xi} & \dotsc & 


\dst\int_{a_1}\dst\frac{\tau^{h-1}d\tau}{\xi} \\ 


\dotsc & \dotsc & \dotsc \\ 


\dst\int_{a_h}\dst\frac{d\tau}{\xi} & \dotsc & 


\dst \int_{a_h}\dst\frac{\tau^{h-1}d\tau}{\xi} 


\end{vmatrix} 


}\qquad 


\begin{matrix} 


(j=\ov{1,h}), \\ (z,w)\in\frak R, \\ (\tau,\xi)\in\frak R, 


\end{matrix} 


$$ 


где $\delta_{jk}$ --- символ Кронекера. 


 


Посчитаем индекс коэффициента задачи (20) ([4], c.\,132). В каждой точке 


$(r_k,0)$, $k=\ov{1,2h+2}$, контура $L$ начинаются и оканчиваются по две 


кривые из множества кривых 


$L_j$. Обозначим через $L_{k,1}$, $L_{k,2}$ те кривые, которые 


оканчиваются в произвольно 


зафиксированной точке $(r_k,0)$, $k=\ov{1,2h+2}$,


а через $L_{k,3}$, $L_{k,4}$ --- те 


кривые, которые начинаются в 


указанной точке. Через $\Delta_{L_{k,j}}(r_k-0)$ ($j=1,2$) обозначим 


соответственно предельные значения 


слева функции $\Delta(x,\vartheta)$ в концевой точке $(r_k,0)$ по 


кривым $L_{k,j}$ ($j=1,2$). Аналогично через $\Delta_{L_{k,j}}(r_k+0)$ 


($j=3,4$) обозначим соответственно предельные значения справа функции 


$\Delta(x,\vartheta)$ в 


начальной точке $(r_k,0)$ 


по кривым $L_{k,j}$ ($j=3,4$). С учетом введенных обозначений индекс 


коэффициента 


задачи (20) будет вычисляться по формуле 


$\varkappa=\sum\limits^{2h+2}_{k=1}[\varkappa_k]$, 


где 


$\varkappa_k=\frac{1}{2\pi}(\arg\Delta_{L_{k,1}}(r_k-0)+ 


\arg\Delta_{L_{k,2}}(r_k-0)-\arg\Delta_{L_{k,3}}(r_k+0)- 


\arg\Delta_{L_{k,4}}(r_k+0))$. 


Здесь и в последующем квадратными скобками будем обозначать целую 


часть. Заметим, что $\varkappa_k=\varkappa_{2h+3-k}$ ($k=\ov{1,h+1}$) в 


силу свойства (19) коэффициента $\Delta(x,\vartheta)$. Введем 


дивизор 


$E=(r_1,0)^{[\varkappa_1]}(r_2,0)^{[\varkappa_2]}\cdot\dotsc\cdot 


(r_{2h+2},0)^{[\varkappa_{2h+2}]}$, $\ord E=\varkappa$. 


Задача (20) сводится к последовательному решению двух задач: 


проблемы обращения Якоби ([4], \S\,\S\,4,~5) и задачи нахождения 


мероморфных всюду на $\frak R$ функций, кратных заданному дивизору. 


 


Решение однородной задачи (20) будем искать в виде ([5], c.\,112) 


\begin{multline} 


\wh F(z,w)=\wh\varphi(z,w)\exp\bigg\{\frac{1}{2\pi i}\int_L 


\ln\Delta(\tau,\xi)\wt\omega_{(z,w)}(\tau,\xi)d\tau- \\ 


% 


-\sum^h_{j=1}\bigg[\fuckprime\int^{(z_j,w_j)}_{(z_0,w_0)} 


\wt\omega_{(z,w)}(\tau,\xi)d\tau-n_j\int_{a_j} 


\wt\omega_{(z,w)}(\tau,\xi)d\tau-m_j\int_{b_j} 


\wt\omega_{(z,w)}(\tau,\xi)d\tau\bigg]\bigg\}, 


\end{multline} 


где $(z_0,w_0)$ --- произвольная фиксированная точка поверхности 


$\frak R$, не принадлежащая 


контуру $L$, а через $\wt\omega_{(z,w)}(\tau,\xi)d\tau$ обозначено 


выражение 


$$ 


\wt\omega_{(z,w)}(\tau,\xi)d\tau= 


\frac{w+\xi}{2\xi}\frac{d\tau}{\tau-z}- 


\frac{\wt w+\xi}{2\xi}\frac{d\tau}{\tau-\wt z}. 


$$ 


Здесь $(\wt z,\wt w)$ --- произвольная фиксированная точка поверхности 


$\frak R$, не 


принадлежащая контуру $L$, причем $(\wt z,\wt w)\ne(z_0,w_0)$. Штрих 


перед интегралом в формуле 


(22) означает, что путь интегрирования не пересекает канонические 


сечения. Точки $(z_j,w_j)$ и целые числа $n_j$, $m_j$ ($j=\ov{1,h}$) 


находятся из того условия, чтобы 


входящая в формулу (22) экспонента не имела существенно особых точек 


при $(z,w)=(\infty,\infty_1)$ 


и $(z,w)=(\infty,\infty_2)$. Разложим $w(z)$ и $\frac{1}{\tau-z}$ в 


ряды по степеням $z$ при $z\to\infty$ и подставим эти 


разложения в формулу (22). Заметим, что 


$$ 


w(z)\sim z^{h+1}+C_h z^h+C_{h-1}z^{h-1}+\dotsb+C_0\quad 


\text{при}\ \; z\to\infty, 


$$ 


где $C_i$ ($i=\ov{0,h}$) --- некоторые 


постоянные. Приравнивая нулю соответствующие коэффициенты при 


одинаковых положительных степенях $z$, получим однородную треугольную 


систему $h$ уравнений с $h$ неизвестными $\lambda_i$ 


\begin{equation} 


\left\{ 


\begin{split} 


\lambda_0&=0,\\ 


\lambda_0C_h+\lambda_1&=0,\\ 


\lambda_2C_{h-1}+\lambda_1C_h+\lambda_2&=0,\\ 


\dotsc\dotsc\dotsc\dotsc&\  \\ 


\lambda_0C_2+\lambda_1C_3+\dotsb+\lambda_{h-2}C_h+ 


\lambda_{h-1}&=0, 


\end{split} 


\right. 


\end{equation} 


где $\lambda_i$ выражаются формулой 


$$ 


\lambda_i=\frac{1}{2\pi i}\int_L 


\frac{\ln\Delta(\tau,\xi)}{2\xi}\tau^id\tau- 


\sum^h_{j=1}\bigg[\fuckprime 


\int^{(z_j,w_j)}_{(z_0,w_0)}\frac{\tau^i}{2\xi}d\tau- 


n_j\int_{a_j}\frac{\tau^i}{2\xi}d\tau- 


m_j\int_{b_j}\frac{\tau^i}{2\xi}d\tau\bigg], 


\quad i=\ov{0,h-1}. 


$$ 


Очевидно, что система (23) имеет лишь тривиальное решение. Таким 


образом, получили задачу обращения абелевых интегралов первого 


рода, которая имеет вид 


\begin{equation} 


\frac{1}{2\pi i}\int_L 


\frac{\ln\Delta(\tau,\xi)}{2\xi}\tau^id\tau=\sum^h_{j=1}\bigg[ 


\fuckprime\int^{(z_j,w_j)}_{(z_0,w_0)} 


\frac{\tau^i}{2\xi}d\tau-n_j\int_{a_j}\frac{\tau^i}{2\xi}d\tau- 


m_j\int_{b_j}\frac{\tau^i}{2\xi}d\tau\bigg],\quad 


i=\ov{0,h-1}. 


\end{equation} 


Эта задача заключается в нахождении точек $(z_j,w_j)$ и целых чисел 


$n_j$, $m_j$ ($j=\ov{1,h}$), 


удовлетворяющих (24). Это и есть проблема обращения Якоби. Общая идея 


способа, c помощью которого получены уравнения (24), изложена в 


работе ([5], c.\,112). Более подробно этот способ применен в статье 


([6], 


c.\,48). Для вычисления точек $(z_j,w_j)$ и целых чисел $n_j$, $m_j$ 


($j=\ov{1,h}$) годится описанный в статье 


[4] алгоритм решения задачи обращения Якоби. В указанной работе 


приведен алгоритм для случая произвольной римановой поверхности рода 


$h$. Специфика гиперэллиптического случая отражена в статье [7], 


поэтому на алгоритме здесь не останавливаемся. Будем считать, что 


решение проблемы обращения Якоби найдено и тем самым определен 


дивизор $F=(z_0,w_0)^h(z_1,w_1)^{-1}(z_2,w_2)^{-1}\cdot\dotsc\cdot 


(z_h,w_h)^{-1}$. 


 


Функция $\wh\varphi(z,w)$, входящая в формулу (22), есть произвольная 


мероморфная всюду 


на $\frak R$ функция, кратная дивизору 


\begin{multline*} 


J^{-1}E^{-1}F^{-1}=(\infty,\infty_1)^1 


(\infty,\infty_2)^1(r_1,0)^{-[\varkappa_1]} 


(r_2,0)^{-[\varkappa_2]}\cdot\dotsc\cdot \\


%


\cdot (r_{2h+2},0)^{-[\varkappa_{2h+2}]}(z_0,w_0)^{-h} 


(z_1,w_1)^1(z_2,w_2)^1\cdot\dotsc\cdot(z_h,w_h)^1. 


\end{multline*}  


Функцию $\wh\varphi(z,w)$ можно искать в виде 


$$ 


\frac{P(z)w+Q(z)}{R(z)},


$$ 


где $P(z)$, $Q(z)$, $R(z)$ --- многочлены от $z$, 


степени и коэффициенты 


которых находятся из 


условий, обеспечивающих кратность. Автором получен конкретный вид 


функции $\wh\varphi(z,w)$, кратной дивизору $J^{-1}E^{-1}F^{-1}$, однако 


на этом здесь не останавливаемся. 


Обозначим через $l$ число произвольных комплексных постоянных, 


содержащихся в функции $\wh\varphi(z,w)$, 


а следовательно, и в общем решении 


однородной задачи (20). Тогда общее решение этой задачи выражается 


формулой (22) и содержит $l$ линейно независимых решений. 


 


Перейдем к решению однородной задачи для дифференциалов (21). Будем 


искать его в виде 


\begin{multline} 


d\Psi(z,w)=d\psi(z,w)\exp\bigg\{-\frac{1}{2\pi i} 


\int_L\ln\Delta(\tau,\xi)\wt\omega_{(z,w)}(\tau,\xi)d\tau+ \\ 


% 


+ \sum^h_{j=1}\bigg[


\fuckprime\int^{(z_j,w_j)}_{(z_0,w_0)}\wt\omega_{(z,w)} 


(\tau,\xi)d\tau 


-n_j\int_{a_j}\wt\omega_{(z,w)}(\tau,\xi)d\tau- 


m_j\int_{b_j}\wt\omega_{(z,w)}(\tau,\xi)d\tau\bigg]\bigg\},


\end{multline} 


где $d\psi(z,w)$ --- произвольный абелев дифференциал, кратный дивизору 


$JEF$, а все 


остальные символы в правой части (25) те же, что и в правой части 


(22). Дифференциал $d\psi(z,w)$ можно искать в виде 


$$ 


\frac{P^*(z)w+Q^*(z)}{R^*(z)}dz, 


$$ 


где $P^*(z)$, $Q^*(z)$, $R^*(z)$ --- многочлены от $z$, степени и 


коэффициенты которых находятся 


из условий, обеспечивающих кратность дивизору 


\begin{multline*}  


JEF=(\infty,\infty_1)^{-1}


(\infty,\infty_2)^{-1}(r_1,0)^{[\varkappa_1]} 


(r_2,0)^{[\varkappa_2]}\cdot\dotsc\cdot \\


%


\cdot(r_{2h+2},0)^{[\varkappa_{2h+2}]}(z_0,w_0)^{h} 


(z_1,w_1)^{-1}(z_2,w_2)^{-1}\cdot\dotsc\cdot(z_h,w_h)^{-1}. 


\end{multline*}  


Вид дифференциала $d\psi(z,w)$


так же, как и вид функции $\wh\varphi(z,w)$, можно 


конкретизировать, но на этом здесь не останавливаемся. Обозначим через 


$l'$ число произвольных 


комплексных постоянных, содержащихся в дифференциале $d\psi(z,w)$, а 


следовательно, и в общем решении задачи (21). Тогда общее решение 


этой задачи выражается формулой (25) и содержит $l'$ линейно независимых 


решений. Таким образом, получены общие решения однородных задач (20) 


и (21) без учета условия симметрии (17). Проведем рассуждения, 


позволяющие учесть это условие. 


\medskip 


 


{\bf 6. Утверждение.} {\it 


Число произвольных комплексных постоянных, 


содержащихся в общем несимметричном решении задачи $(20)$, равно числу 


произвольных вещественных постоянных, содержащихся в общем 


симметричном решении этой задачи. 


} 


\medskip 


 


\begin{pf} 


Заметим, что если $F(z,w)$ есть решение задачи (20), то 


$\ov{F(-\ov z,-w(\ov z))}$ 


также 


есть решение этой задачи в силу (19). Следовательно, 


$\frac{F(z,w)+\ov{F(-\ov z,-w(\ov z))}}{2}$, 


$\frac{F(z,w)-\ov{F(-\ov z,-w(\ov z))}}{2i}$ 


также являются 


решениями задачи (20), причем они удовлетворяют условию симметрии 


(17). Кроме того, любое решение $F(z,w)$ однородной задачи (20) может 


быть представлено в виде линейной комбинации симметричных решений, 


указанных выше. Пусть $\{F_k(z,w)\}$, $k=\ov{1,l}$, является базисом 


пространства несимметричных 


решений задачи (20). Тогда система симметричных решений 


$\frac{F_k(z,w)+\ov{F_k(-\ov z,-w(\ov z))}}{2}$, 


$\frac{F_k(z,w)-\ov{F_k(-\ov z,-w(\ov z))}}{2i}$ 


представляет 


собой полную систему решений однородной задачи (20). Эта система 


содержит $2l$ функций. Выберем в ней $l$ линейно независимых симметричных 


функций и разложим решение однородной задачи по новому базису. 


Получим формулу общего симметричного решения задачи (20) тогда и 


только тогда, когда все коэффициенты этого разложения будут 


вещественными. Количество произвольных вещественных постоянных, 


содержащихся в общем симметричном решении задачи (20), будет равно $l$. 


\end{pf} 


 


Из утверждения следует, что картина разрешимости однородной задачи 


(20) останется той же самой и с учетом условия симметрии (17). Под 


картиной разрешимости понимается количество линейно независимых 


решений (соответственно и произвольных постоянных), содержащихся в 


общем решении задачи. В случае симметричного решения речь идет о 


произвольных вещественных постоянных. В доказательстве утверждения 


фактически описана процедура нахождения общего решения однородной 


задачи (20), удовлетворяющего условию симметрии (17). Будем считать, 


что оно построено, и обозначим его $\wh F_*(z,w)$. 


\medskip 


 


{\bf 7.} Теперь рассмотрим неоднородную задачу (18). Для построения ее 


симметричного общего решения достаточно построить симметричное частное решение и прибавить к нему общее решение однородной задачи 


$\wh F_*(z,w)$. Построим сначала частное решение, не удовлетворяющее 


условию 


симметрии (17), а затем получим симметричное частное решение. 


 


Для разрешимости неоднородной задачи необходимо и достаточно, чтобы 


выполнялись равенства ([4], с.\,146) 


\begin{equation} 


\int_L G_1(\tau,\xi)d\Psi^+(\tau,\xi)=0 


\end{equation} 


для всех линейно независимых решений, содержащихся в $d\Psi(z,w)$. 


Предположим, что 


эти равенства выполняются, и построим частное решение неоднородной 


задачи (18). Частное решение этой задачи будет представлено 


различными аналитическими выражениями в зависимости от того, 


разрешима нетривиально задача (21) или нет. 


 


Рассмотрим первый случай. 


Пусть задача (21) имеет нетривиальное решение. Обозначим его 


$d\Psi_0(z,w)$. Тогда 


задачу (18) можно переписать в виде следующей задачи ``о скачке'' для 


нахождения кусочно-аналитического дифференциала $F(z,w)d\Psi_0(z,w)$: 


$$ 


F^+(x,\vartheta)d\Psi^+_0(x,\vartheta)- 


F^-(x,\vartheta)d\Psi^-_0(x,\vartheta)= 


G_1(x,\vartheta)d\Psi^+_0(x,\vartheta),\quad 


(x,\vartheta)\in L. 


$$ 


В силу (26) эта задача разрешима, а алгоритм ее решения подробно 


описан в ([4], c.\,147--148). Решив эту задачу, легко найти частное 


решение задачи (18). 


 


Рассмотрим второй случай. 


Пусть задача (21) не имеет нетривиальных решений. Это означает, что 


размерность линейного пространства абелевых на $\frak R$ дифференциалов, 


кратных дивизору $JEF$, равна нулю. Отсюда следует


$\ord(JEF)\ge h-1$. Добавляя в 


знаменатель дивизора $JEF$ конечное число точек, не лежащих на $L$, 


можем 


дополнить его до минимального дивизора ([4], с.\,122). Напомним, что 


дивизор $\Delta_{\min}$ называется минимальным, если не существует 


мероморфных всюду 


на $\frak R$ функций, кратных $\Delta^{-1}_{\min}$, и не существует 


абелевых на $\frak R$ дифференциалов, 


кратных $\Delta_{\min}$. Из теоремы Римана--Роха ([8], c.\,122) следует, 


что для минимального 


дивизора имеет место равенство $\ord\Delta_{\min}=h-1$. Выбор точек, 


которые нужно добавить 


в знаменатель дивизора $JEF$, чтобы дополнить его до минимального, 


основывается на одном из свойств сопряженных пар точек ([4], c.\,157) 


гиперэллиптической поверхности $\frak R$ (т.\,е. пар вида $(z_0,w_0)$, 


$(z_0,-w_0)$ точек указанной 


поверхности). Точки $(\infty,\infty_1)$, $(\infty,\infty_2)$ образуют 


сопряженную пару. Каждая точка 


ветвления, взятая с кратностью два, также порождает сопряженную пару. 


Известно, что сопряженные пары (и только они) обладают тем свойством, 


что существуют две линейно независимые функции 


$$ 


C_1+\frac{C_2}{z-z_0},\quad C_1+\frac{C_2}{z-r_k},\quad 


C_1+C_2z, 


$$ 


кратные соответственно дивизорам 


$$ 


(z_0,w(z_0))^{-1}(z_0,\,-w(z_0))^{-1};\quad 


(r_k,0)^{-2};\quad(\infty,\infty_1)^{-1} 


(\infty,\infty_2)^{-1}. 


$$ 


Если же точки $(z_1,w(z_1))$, $(z_2,w(z_2))$ не образуют сопряженной 


пары, то существует только одна 


функция (постоянная), кратная дивизору 


$(z_1,w(z_1))^{-1}$ $(z_2,w(z_2))^{-1}$. 


Принимая во внимание  


изложенное выше свойство, выберем точку поверхности $\frak R$, не 


лежащую на 


контуре $L$, так, чтобы она не образовывала сопряженных пар с точками 


дивизора $JEF$. Добавим в знаменатель дивизора $JEF$ эту точку, 


взятую с 


кратностью $\varkappa-h-1$. Обозначим полученный дивизор $\Delta_{\min}$. 


Его порядок будет равным $h-1$. Дивизор $\Delta_{\min}$


мы построили таким образом, что 


не существует 


абелевых на $\frak R$ дифференциалов, кратных $\Delta_{\min}$, и не 


существует 


мероморфных всюду на $\frak R$ функций, кратных $\Delta^{-1}_{\min}$. 


Следовательно, $\Delta_{\min}$ и есть 


минимальный дивизор. Принимая его за характеристический дивизор 


мероморфного аналога ядра Коши ([4], c.\,124), построим последний по 


формуле из ([4], c.\,125). Обозначим его $d\omega_{(z,w)}(\tau,\xi)$. 


Через $X_0(z,w)$ обозначим функцию, в 


которую переходит правая часть формулы (22), если в ней положить 


$\wh\varphi(z,w)\equiv1$. 


Функция $X_0(z,w)$ удовлетворяет краевому условию (20), поэтому задачу 


(18) 


можно свести к задаче ``о скачке''


$$ 


\frac{F^+(x,\vartheta)}{X^+_0(x,\vartheta)}- 


\frac{F^-(x,\vartheta)}{X^-_0(x,\vartheta)}= 


\frac{G_1(x,\vartheta)}{X^+_0(x,\vartheta)}, 


\quad (x,\vartheta)\in L. 


$$ 


Применяя формулу для решения этой задачи ([4], с.\,149), 


получим частное несимметричное 


решение 


$$ 


F(z,w)=\frac{X_0(z,w)}{2\pi i}\int_L 


\frac{G_1(\tau,\xi)}{X^+_0(\tau,\xi)}d\omega_{(z,w)}(\tau,\xi) 


$$ 


задачи (18). 


 


Частное симметричное решение можно 


получить из несимметричного следующим образом: 


\begin{equation} 


F_*(z,w)= 


\frac{F(z,w)+\ov{F(-\ov z,-w(\ov z))}}{2}. 


\end{equation} 


Если $F(z,w)$ --- решение неоднородной задачи, то и 


$\ov{F(-\ov z,-w(\ov z))}$ также является решением этой задачи в силу 


(19). Тогда общее 


симметричное решение задачи (18) будет иметь вид 


\begin{equation} 


F^*(z,w)=\wh F_*(z,w)+F_*(z,w), 


\end{equation} 


где $\wh F_*(z,w)$ --- общее симметричное решение однородной задачи 


(20), а $F_*(z,w)$ --- частное 


симметричное решение задачи (18), вычисляемое по формуле (27). 


\medskip 


 


{\bf 8.} Подытожим все полученные результаты и опишем тем самым картину 


разрешимости краевой задачи (5). Напомним, что через $l$ обозначается 


число линейно независимых симметричных решений однородной задачи (20) 


и число произвольных вещественных постоянных в симметричном общем 


решении неоднородной задачи (18). Если $l=0$, то однородная задача (20) 


имеет лишь тривиальное решение. Через $l'$ обозначается число условий 


разрешимости неоднородной задачи (18). Эти условия описываются 


формулой (26). 


 


\begin{thm} 


Если $\varkappa>2h$, то $l=\varkappa-h-1$, $l'=0$. Если $\varkappa<2$, 


то $l=0$, $l'=-\varkappa+h+1$. В ``особом случае'' $2\le\varkappa\le2h$ 


имеет место точная оценка $\max\{0,\varkappa-h-1\}\le 


l\le\frac{\varkappa}{2}$, $l'=l-\varkappa+h+1$. Во всех случаях, когда 


решение задачи $(18)$ 


существует, оно вычисляется по формуле $(28)$. 


\end{thm} 


 


Перейдем теперь от решения задачи (18) к решению задачи (5) . Так как 


построенное решение задачи (18) удовлетворяет условию симметрии (17), 


то подставив вместо $w$ в решение $F^*(z,w(z))$ выражение 


$\sqrt{\prod\limits^{2h+2}_{k=1}(z-r_k)}$, получим тем самым 


решение $\Phi(z)$ задачи (5). Выполнив обратное преобразование Фурье, 


найдем 


функцию $\varphi(x)$. Она является решением уравнения (1), ядра которого 


подчиняются условиям (4), (7)--(10). 


 


Сформулируем окончательный результат. 


 


\begin{thm} 


Уравнение $(1)$, ядра и свободный член которого подчиняются 


условиям $(4)$, $(7)$--$(10)$, допускает решение в замкнутой форме. Это 


уравнение эквивалентно краевой задаче Римана $(5)$, картина 


разрешимости которой полностью определена в теореме~$1$. Во всех 


случаях, когда решение уравнения $(1)$ существует, оно вычисляется по 


формуле 


$$


\varphi(x)=\frac{1}{\sqrt{2\pi}} 


\int^\infty_{-\infty}F^*\bigg(t, 


\sqrt{\prod^{2h+2}_{k=1}(t-r_k)}\,\bigg) 


e^{-itx}dt, 


$$


где $F^*(z,w)$ определяется формулой $(28)$. 


\end{thm} 


 


\begin{thebibliography}{99} 


 


\bibitem{1} 


Комяк И.И. {\em Об интегральных уравнениях типа свертки с сопряжением} 


// Весцi АН БССР. Сер. фiз.-матэм. наук. -- 1967. -- \No\,2. --  


С.\,33--47. 


\bibitem{2} 


Гахов Ф.Д., Черский Ю.И. {\em Уравнения типа свертки}. -- M.: Наука, 


1978. -- 296 c. 


\bibitem{3} 


Евграфов М.А. {\em Аналитические функции}. -- М.: Наука, 1991. -- 447 c. 


\bibitem{4} 


Зверович Э.И. {\em Краевые задачи теории аналитических функций в 


г\"ельдеровских 


классах на римановых поверхностях} // УМН. -- 1971. -- 


Т.ХХVI. --  


Вып.\,1. -- С.\,113--179. 


\bibitem{5} 


Зверович Э.И. {\em Смешанная задача теории упругости для плоскости


с разрезами, лежащими на вещественной оси} //


Тр. симпозиума по механ. сплошн. среды и 


родственным 


пробл. анализа. -- Тбилиси, 1971. -- T.\,1. -- C.\,103--114. 


\bibitem{6} 


Корзан Л.А. {\em Однородная смешанная задача теории упругости для 


полосы с разрезами, лежащими на вещественной оси} // Весцi АНБ. 


Сер. фiз.-матэм. наук. -- 


1996. -- \No\,4. -- С.\,44--49. 


\bibitem{7} 


Зверович Э.И. {\em О конструктивном решении краевой задачи Римана на 


гиперэллиптических римановых поверхностях} // ДАН СССР. -- 1971. -- 


T.\,199. --  


\No\,4. -- С.\,758--761. 


\bibitem{8} 


Спрингер Дж. {\em Введение в теорию римановых поверхностей}. --М.: Ин. 


лит., 1960. -- 343 c. 


 


\end{thebibliography} 


 


\vskip 2cm 


 


\post{\it Белорусский Государственный}{\it Поступила} 


\post{\it Технологический Университет}{23.03.1998} 


 


\label{end} 


\end{document} 


